\section{Coverage of the Target Manifold, and Why AUC Is Not Used}
\label{app:coverage}

Section~\ref{sec:why} measures how much of the real target manifold a set covers as $k$-NN recall
at $k = 5$~\citep{kynkaanniemi2019improved}: the fraction of the $4040$ target images that fall inside
the $k$-NN balls of the compared set, in each of the three embedding spaces of Table~\ref{tab:silhouette}.
Every real set covers the target near its own ceiling; both synthetic pools collapse, most sharply in
\texttt{clipL224}.

\begin{center}\small
\begin{tabular}{lccc}
\toprule
Compared set & \texttt{dinoL224} & \texttt{dinoL518} & \texttt{clipL224} \\
\midrule
target against itself (random halves) & 0.9332 & 0.9431 & 0.8931 \\
COD10K part vs CAMO part of the target & 0.8898 & 0.9003 & 0.8730 \\
NC4K, a different real set & 0.8874 & 0.8995 & 0.7921 \\
raw HKU-IS, the generator's input & 0.7473 & 0.7097 & 0.7470 \\
authors' synthetic pool & 0.4829 & 0.4834 & 0.1277 \\
our synthetic renders & 0.4668 & 0.5421 & 0.1349 \\
\bottomrule
\end{tabular}
\end{center}

A real-versus-synthetic classifier is not used as evidence, because it cannot distinguish a
distributional gap from a re-encoding artifact: re-saving the \emph{identical} target images at
JPEG quality $30$ already separates them from themselves at AUC $0.9928$ in \texttt{clipL224}, against
$0.9993$ and $0.9995$ for the two synthetic pools. \emph{Source:} \texttt{rebuild/A3/out/a3\_coverage.csv},
\texttt{a3\_probe\_table.csv}.
